\section{Related work}
\label{sec:related}

\paragraph{Spectral theory.}
The standard account of moir\'e is spectral. Superposing two transparencies
multiplies their intensities, so their spectra convolve, and the convolution
produces sum and difference components that were in neither layer. When two
carriers are close, the $(1,-1)$ difference component falls below the eye's cutoff
and appears as a large, slow pattern. Amidror's monograph is the complete
treatment~\cite{Amidror2009Theory}, developed with Hersch across a sequence of
papers that moves from analysis of screen interactions~\cite{Amidror1994Spectral},
to phase shifts~\cite{Amidror1996Fourier}, to synthesis under geometric
transformation~\cite{Amidror1998Fourier}, and finally to aperiodic
layers~\cite{Amidror2003Moire, AmidrorII}, whose random-dot extreme is the Glass
pattern~\cite{Glass1969Moire}. The account is still growing: Feng extends it to
\emph{objective} structures --- ring, helical and lattice patterns whose members
are related by progressive symmetry transformations --- through an augmented
Fourier approach~\cite{Feng2025Objective}. Those structures are the closest
formal relatives of our walking families (\S\ref{sec:walking}), reached from the
frequency side; we reach them per pixel, with a certificate, because a renderer
has to. That is the general shape of our debt to this theory: it tells you
\emph{that} a fringe exists and where its frequency lies, but it is a statement
about spectra of images, not a renderer.

\paragraph{The indicial equation.}
Running alongside the spectral account is an older and more elementary one. Label
the members of each layer, ask where a member of one crosses a member of the other
in a fixed relation, and solve: Oster, Wasserman and Zwerling derive the fringe
curves of the standard configurations this way~\cite{Oster1964}, Amidror gives
it as the geometric approach, complementary to the spectral one --- curved
gratings read as indexed line families (Vol.~I, \S11.2
of~\cite{Amidror2009Theory}) --- and the metrology literature develops it at
book length~\cite{Bryngdahl1974Moire, Patorski1993Handbook}. This is the direct ancestor of our
Section~\ref{sec:fringe}, and our index field is its label made continuous and
evaluated per pixel. We claim no credit for the locus. What the indicial method
does not supply is the quantity a renderer needs --- the coverage \emph{between} the
fringes --- nor any statement of when the construction is meaningful; our
Theorem~\ref{thm:fringe} supplies the first as a closed-form profile of the actual
compositing operator and \S\ref{sec:visibility} the second as a field that can be
drawn. Both are per-pixel, exact rather than asymptotic, and indifferent to whether
anything is periodic.

\paragraph{Constructive synthesis.}
Band moir\'e~\cite{Hersch2004Band} is the landmark generation method in graphics:
a base layer of image bands, replicated and compressed along one axis, sampled by a
revealing line grating, with a derived linear map from band space to moir\'e space
that predicts the revealed image's orientation, scale and velocity. Applying
nonlinear coordinate transforms to both layers yields curved layers whose
superposition is still a clean periodic moir\'e. Level-line
moir\'e~\cite{Chosson2014Beating} changes the encoding: a bilevel target becomes an
elevation profile, embedded as small local displacements of the base grating, and
the superposition reveals the profile's \emph{level lines}, which traverse the
shape as the layers slide. Walger and Hersch multiplex up to seven such images in
one base~\cite{Walger2015Hiding}. Target-driven phase
modulation~\cite{Tsai2013Target} is the strongest inverse formulation: given a
grayscale target, solve for two curvilinear gratings whose dominant difference
component reproduces it, with a smoothness term distributing information between
the layers and an appearance phase to disguise them. Lebanon and
Bruckstein~\cite{Lebanon2001Variational} give the variational counterpart.

These methods and ours are close in spirit and disjoint in mechanism. They solve
an inverse problem to \emph{produce} two layers; we make the forward problem cheap
enough that a person can solve the inverse problem by hand, in real time, at any
zoom. The level-line formulation is the closest relative, and the relationship is
clean: their construction embeds a designed elevation profile in a grating's phase
so that its contours appear as fringes, and our Section~\ref{sec:beyond} shows
that in the field model this is not a construction at all but an immediate
consequence of the fringe law, available for any field, on screen, with no printing
step. We regard that as evidence the representation is the right one, not as a
claim to have invented level-line moir\'e.

\paragraph{Halftoning.}
Halftoning is where moir\'e is classically a hazard and occasionally a tool.
Rotated dispersed dither disperses the spectral impulses of a Bayer array so that
colour separations do not beat~\cite{Ostromoukhov1994Rotated, Ulichney1987};
multi-colour dithering extends the control to non-standard
inks~\cite{Ostromoukhov1999MultiColor}. Artistic
screening~\cite{Ostromoukhov1995Artistic} turns the halftone dot into an
artist-supplied contour that grows with intensity and can be pushed through a
conformal map, which makes the screen itself the designed object --- the direct
conceptual ancestor of band moir\'e, and a useful reminder that the expressive
material here is the \emph{carrier}, not the picture.

\paragraph{Physical realisation and sensing.}
Moir\'e magnifiers pair a micro-image array with a slightly mismatched lens
array~\cite{Hutley1994Moire}. EPFL's later work replaces printed revealers with
cylindrical microlenses~\cite{Walger2019Moire, Walger2020LevelLine} and Saberpour
et al.\ carry designed moir\'e onto curved fabricated
surfaces~\cite{Saberpour2020Fabrication}. Refractive
steganography~\cite{Papas2012Magic} and metallic-substrate
hiding~\cite{Pjanic2015Color} pursue the same hidden-image goal by other optics.
Running the amplification backwards turns a moir\'e into a precise passive
encoder: MoiréBoard for camera pose~\cite{Xiao2021MoireBoard}, MoiréTag for
angle~\cite{Qiu2023MoireTag}, MoiréWidgets for tangible
controls~\cite{CamposZamora2024MoireWidgets}. These works establish that the phase
of a moir\'e is a quantity worth measuring to high precision, which is indirectly
an argument for having it available in closed form.

\paragraph{Interactive authoring.}
FabObscura~\cite{Sethapakdi2025FabObscura} is the nearest system-shaped relative:
it parameterises barrier-grid occlusion patterns as functions, and couples a
pattern editor to a live canvas and a fabrication export. Barrier-grid animation
is a parallax technique rather than interference, so the mechanism differs, but its
design stance --- find the right parameterisation, then expose it directly rather
than making the user reason about pixels --- is the one we tried to adopt. Outside
archival venues, procedural moir\'e is a staple of shader practice, going back to
Glassner's tutorial treatment~\cite{Glassner1997Inside} and Sen's product-delay
constructions~\cite{Sen1998Product, Sen2000Moire}, in the wider tradition of
evaluating a pattern per point rather than storing it~\cite{Ebert2003Texturing}.
The shader-sandbox idiom is
essentially universal and essentially always the same: evaluate two periodic
functions and combine them, typically as \texttt{abs(p1 - p2)}. That is a
\emph{value} composite rather than a distance composite, and
Section~\ref{sec:zoom} shows what it costs --- there is no stroke width to
stabilise, so the pattern is a different pattern at every zoom level, and the
displayed moir\'e is partly the sampler's own.

\paragraph{Moir\'e in visualization.}
Direct use of interference to encode data is rare. MoireGraphs is the notable
appearance of the word in the visualization literature, and it names a
resemblance: a radial focus+context layout whose concentric rings of visual nodes
\emph{look} like a circular moir\'e~\cite{JankunKelly2003Moire}. The place where
fringes genuinely carry data is experimental mechanics and optical metrology.
Moir\'e interferometry reads in-plane displacement: fringes are contour maps of
displacement, with fringe order $N_x=fU$ for reference grating frequency $f$,
and strains follow by differencing fringe orders~\cite{Post2008Moire,
Post1994High}. Moir\'e topography reads surface height the same way --- its
fringes literally are elevation contours~\cite{Takasaki1970Moire,
Meadows1970Contours} --- and Takeda's Fourier-transform method turned fringe
images back into phase maps numerically~\cite{Takeda1982Fourier}. These are
measurement techniques, and their governing relation is exactly our fringe law
read in the other direction. Section~\ref{sec:beyond} takes it as a design
primitive instead: the encoded field is arbitrary, on screen, and live.

\paragraph{Distance fields and certified search.}
The machinery we borrow from is the distance-field tradition: guaranteed
intersection from a Lipschitz constant~\cite{Kalra1989}, sphere tracing, which uses
such a bound to skip space it can prove is empty~\cite{Hart1996Sphere}, and its
sharpening by \emph{local} Lipschitz bounds along a segment~\cite{Galin2020};
adaptively sampled distance fields as a shape
representation~\cite{Frisken2000Adaptively}; and the now-standard practice of
shading vector art from a distance rather than a coverage
mask~\cite{Green2007Improved, Quilez, Quilez2019Distance}. Our index window
(Section~\ref{sec:window}) is a sphere-tracing argument moved from space into
\emph{index} space, and our skip rule is segment tracing along the index axis:
where a ray marcher steps by the distance it can prove is empty, we step by the
number of \emph{rings} we can prove cannot be nearest. The alternative discipline,
enumerating every cell a ray meets~\cite{Amanatides1987}, is what the fixed sweep
we replace amounts to. The closest methodological neighbours are the range-analysis
renderers, which bound a function over a spatial cell and subdivide until the
bound decides~\cite{Knoll2009, Keeter2020, Sharp2022}. Section~\ref{sec:discussion}
argues that our domain being $\Z$ rather than $\R^2$ is what upgrades their
tolerance into a certificate. Seen from optimization, the certified scan is a Piyavskii--Shubert
search~\cite{Piyavskii1972, Shubert1972} run over $\Z$ instead of an interval of
$\R$: the same Lipschitz envelopes, with the stopping tolerance replaced by
integer exhaustion. The support-function facts we need are standard
convex geometry~\cite{Rockafellar1970Convex, Schneider2014Convex}, and the same
object is what collision detection queries when it walks a Minkowski
difference~\cite{GJK1988, Ericson2004} --- though there it prices one pair of
bodies, where here the certificate ranges over an indexed family of them; our
contribution is noticing which four of its properties this problem is asking
for.

\paragraph{Prefiltering a pattern.}
Where a procedural pattern is too fine to sample, the correct answer is to
integrate it rather than to point-sample harder, a line that runs from the
frequency analysis of aliasing~\cite{Crow1977, Mitchell1990} through the
analytically prefiltered procedural textures surveyed by Lagae et
al.~\cite{Lagae2010} to normal-distribution methods for glinty
surfaces~\cite{Yan2014, Yan2016}. Two things connect that work to ours. Our
saturated regime (\S\ref{sec:limits}) is exactly a place where a correct analytic
mean would replace a fallback we know to be approximate, and we do not have one for
a walking family. And the glint problem has our structure --- ``which member of a
dense indexed family lies within this footprint'' --- solved with a
budget~\cite{Yan2016} where an index interval might certify.
